\documentclass[10pt,a4paper]{amsart}
\usepackage[margin=19mm]{geometry}
\usepackage{amsmath,amssymb,mathtools}
\usepackage[scr=rsfs,cal=euler]{mathalfa}
\usepackage{xfrac}
\usepackage[unicode,hidelinks,bookmarks=false]{hyperref}
\newcommand{\ie}{i.e.\ }
\newcommand{\cf}{cf.\ }
\newcommand{\resp}{respectively\ }
\newcommand{\Subsection}{\subsection}
\providecommand{\nobreakdash}{\nobreak\hbox{-}}
\setlength{\emergencystretch}{2em}
\multlinegap0pt
\usepackage{amssymb}
\usepackage{algorithm}\usepackage{algpseudocode}
\usepackage{tikz}
\newtheorem{lemma}{Lemma}
\title{Statistical Properties of Nonparametric MLE under Laplace Noise}
\author{Yifei Xiong, Nianqiao Phyllis Ju, Vinayak Rao}
\date{Source: arXiv:2608.25997v1 (CC BY 4.0)}
\begin{document}
\maketitle

\noindent\emph{Source and licence.} Statistical Properties of Nonparametric MLE under Laplace Noise. Version arXiv:2608.25997v1 (CC BY 4.0), distributed by arXiv under the Creative Commons Attribution 4.0 International licence, http://creativecommons.org/licenses/by/4.0/, as recorded in the arXiv licence metadata for this exact version and shown on the abstract page https://arxiv.org/abs/2608.25997v1. Full licence notice: \texttt{licenses/arxiv-2608.25997-CC-BY-4.0.txt}.\par\smallskip

\noindent\textit{Editorial scope.} Counted unit: the article's lemma lem:laplace-full-rank (source0.tex lines 857--867). The article's complete original proof of it is delivered with it (source0.tex lines 869--938, one \texttt{proof} environment of 70 lines). No mathematics is written, repaired or re-derived: the statement and the proof are the article's own lines, in the article's own notation, and no retained line is substituted or re-flowed.  No further source-local context is needed: every cross-reference of every retained line resolves inside the two delivered spans.  Nothing was omitted from any retained span, so the package carries no omission record.  No retained line contains a citation, so the wrapper carries no bibliography and no bibliography entry.  No retained line contains a figure environment, an image-inclusion command or drawing code, so no image is delivered and the package declares no asset dependency.  Editorial formatting only, in addition: an AMS \texttt{amsart} preamble that restates the article's own declarations and macros used by the retained lines, namely the environments \texttt{lemma} on separate counters, together with the standard packages those lines need.  The retained environments print in this document as Lemma 1 in order of appearance, whereas the article prints the same items with the numbering of its own full document.  The complete native source of the article as distributed in the arXiv e-print for this exact version is bundled unchanged as \texttt{sources/208612/source0.tex}.
\begin{lemma}[Full column rank of the Laplace likelihood matrix]\label{lem:laplace-full-rank}
Let $u_1<\cdots<u_\nu$ be the distinct values among
\begin{equation*}
\Pi_{[-a,a]}(X_1),\dots,\Pi_{[-a,a]}(X_n),
\end{equation*}
and define the $n\times \nu$ matrix
\begin{equation*}
\mathbf S=(s_{ij}),\quad s_{ij}:=f_b(X_i-u_j)=\frac{1}{2b}\exp\left(-\frac{|X_i-u_j|}{b}\right).
\end{equation*}
Then the columns of $\, \mathbf S$ are linearly independent.
\end{lemma}

\begin{proof}[Proof of Lemma~\ref{lem:laplace-full-rank}]
Suppose that
\begin{equation*}
\sum_{j=1}^\nu c_j s_{ij}=0,\quad i=1,\dots,n,
\end{equation*}
for some coefficients $c_1,\dots,c_\nu$. Define
\begin{equation*}
f(x):=\sum_{j=1}^\nu \frac{c_j}{2b}\exp\left(-\frac{|x-u_j|}{b}\right).
\end{equation*}
Then $f(X_i)=0$ for all $i=1,\dots,n$.
Let $X_{(1)}\le \cdots \le X_{(n)}$ denote the ordered observations. Since $X_{(1)}\le u_1$, evaluating at $X_{(1)}$ gives
\begin{equation}\label{eq:left-tail-identity}
\sum_{j=1}^\nu c_j e^{-u_j/b}=0.
\end{equation}

We now eliminate the coefficients successively. For $k=1,\dots,\nu-2$, the point $u_{k+1}$ lies in $(-a,a)$, so it must equal one of the observations. Hence there exists an index $t_k$ such that
\begin{equation*}
X_{t_k}=u_{k+1}\in (u_k,u_{k+1}].
\end{equation*}
Assume inductively that $c_1=\cdots=c_{k-1}=0$. Then \eqref{eq:left-tail-identity} reduces to
\begin{equation*}
\sum_{j=k}^\nu c_j e^{-u_j/b}=0.
\end{equation*}
Evaluating $f$ at $X_{t_k}=u_{k+1}$ yields
\begin{align*}
0&=\sum_{j=k}^\nu \frac{c_j}{2b}\exp\left(-\frac{|u_{k+1}-u_j|}{b}\right) \\
&=\frac{c_k}{2b}\exp\left(\frac{u_k-u_{k+1}}{b}\right)+\frac{1}{2b}\exp\left(\frac{u_{k+1}}{b}\right)\sum_{j=k+1}^\nu c_j e^{-u_j/b} \\
&=\frac{c_k}{2b}\exp\left(\frac{u_k-u_{k+1}}{b}\right)-\frac{c_k}{2b}\exp\left(\frac{u_{k+1}-u_k}{b}\right) \\
&=-\frac{c_k}{b}\sinh\left(\frac{u_{k+1}-u_k}{b}\right).
\end{align*}
Since $u_{k+1}>u_k$, this implies $c_k=0$. Thus
\begin{equation*}
c_1=\cdots=c_{\nu-2}=0.
\end{equation*}

It remains to treat the last two coefficients $c_{\nu-1}$ and $c_\nu$.

\textbf{Case (i).} If $u_\nu<a$, or if $u_\nu=a$ and some observation equals $a$, then there exists an observation in $(u_{\nu-1},u_\nu]$. Repeating the same argument once more yields $c_{\nu-1}=0$, and then~\eqref{eq:left-tail-identity} gives $c_\nu=0$.

\textbf{Case (ii).} If $u_\nu=a$ and no observation lies in $(u_{\nu-1},a]$. Then there must exist an observation $X_r>a$, since otherwise the projected value $a$ would not occur among
\begin{equation*}
\Pi_{[-a,a]}(X_1),\dots,\Pi_{[-a,a]}(X_n).
\end{equation*}
Using $c_1=\cdots=c_{\nu-2}=0$, equation \eqref{eq:left-tail-identity} becomes
\begin{equation}\label{eq:last-two-left}
c_{\nu-1}e^{-u_{\nu-1}/b}+c_\nu e^{-a/b}=0.
\end{equation}
On the other hand, evaluating $f(X_r)=0$ and using $X_r>a=u_\nu>u_{\nu-1}$ gives
\begin{equation}\label{eq:last-two-right}
c_{\nu-1}e^{u_{\nu-1}/b}+c_\nu e^{a/b}=0.
\end{equation}
From \eqref{eq:last-two-left}, we have
\begin{equation*}
c_\nu=-c_{\nu-1}\exp\left(\frac{a-u_{\nu-1}}{b}\right).
\end{equation*}
Substituting this into \eqref{eq:last-two-right} gives
\begin{align*}
0&=c_{\nu-1}e^{u_{\nu-1}/b}+c_\nu e^{a/b} \\
&=c_{\nu-1}e^{u_{\nu-1}/b}-c_{\nu-1}\exp\left(\frac{a-u_{\nu-1}}{b}\right)e^{a/b} \\
&=c_{\nu-1}\left(e^{u_{\nu-1}/b}-e^{(2a-u_{\nu-1})/b}\right) \\
&=-2c_{\nu-1}e^{a/b}\sinh\left(\frac{a-u_{\nu-1}}{b}\right).
\end{align*}
Since $a>u_{\nu-1}$, we have
\begin{equation*}
\sinh\left(\frac{a-u_{\nu-1}}{b}\right)>0,
\end{equation*}
and therefore $c_{\nu-1}=0$. Then \eqref{eq:last-two-left} yields $c_\nu=0$.

Therefore, $c_1=\cdots=c_\nu=0$, and the columns of $\mathbf S$ are linearly independent.
\end{proof}

\end{document}
